Se logró cumplir el objetivo de construir un contenedor de basura inteligente que detecte entre dos desechos solidos: cartón y plástico, esto para entornos controlados. Teniendo esto en cuenta se puede concluir lo siguiente:

\begin{itemize}
  \item \textbf{Automatización eficiente:} El prototipo desarrollado demuestra la viabilidad de implementar sistemas automatizados basados en Arduino y ESP32 para la clasificación de residuos sólidos, promoviendo la correcta separación de materiales reciclables como plástico y cartón.
  \item \textbf{Contribución ambiental:} La implementación de esta solución tiene el potencial de reducir la contaminación al fomentar la separación de residuos en la fuente, lo que facilita el reciclaje y reduce el impacto ambiental.
  \item \textbf{Uso de tecnologías modernas:} La integración de tecnologías como Edge Impulse para el reconocimiento de imágenes y sensores ultrasónicos mejora la precisión y funcionalidad del sistema, abriendo oportunidades para su expansión a otros tipos de residuos y entornos.
  \item \textbf{Impacto educativo y social:} Este proyecto, aplicado en un entorno académico, también tiene un impacto educativo al sensibilizar a estudiantes y personal sobre la importancia de una correcta gestión de residuos.
  \item \textbf{Limitaciones superables:} Aunque el sistema actual está limitado a clasificar dos materiales, el diseño modular permite una fácil adaptación para manejar más tipos de residuos en el futuro.
\end{itemize}